\section{Discussion}
\label{sec:discussion}
\subsection{Answers to the Research Questions}
\textit{RQ1: checkpoint policy and held-out group.} The fold-specific analysis shows that a pooled near-perfect selected score can obscure a materially weaker fixed-final result on one group allocation. Reporting F1 with class-index MAE and QWK makes the type of error more interpretable. The group audit also resolves a reporting ambiguity: recorded session membership is disjoint, whereas the historical identity warning concerns visual resemblance between two operator-confirmed different tyres. These are different issues and should not be conflated. Checkpoint selection still uses the evaluated validation fold.

\textit{RQ2: localisation as an intervention.} Region agreement does not ensure a downstream gain for a classifier trained on full frames. The fold-specific intervention table localises the overall negative mean: fold 1 worsens, fold 0 is unchanged, and fold 2 has small improvements. This is stronger evidence than either assuming crops help or claiming crops universally harm. It identifies a limitation of this frozen-classifier intervention. A crop-trained classifier or learned fusion could behave differently and requires its own controlled evaluation. Boundary F1 and IoU describe complementary spatial properties, consistent with prior boundary-sensitive evaluation work.

\textit{RQ3: boundary representation.} Point learning improves the average matched error, while the image-level and location-level analyses show nonuniform benefit. The left-middle location changes little and two image averages favour mask-derived boundaries. Both models remain potentially useful: dense masks provide a search region, while point predictions provide direct coordinate proposals. The study compares complete supervised pipelines; it does not identify the backbone, loss, or pretraining source as the cause of the difference.

These results support an evaluation practice: preserve the reference prediction, pair each intervention with it, report the actual decision rule and coverage, and retain invalid cases. The new gate analysis further shows that explanation-guided model eligibility is sensitive to a stricter heuristic threshold. This strengthens the methodological account without turning an exploratory shortlist into a robust selection guarantee.

\subsection{Scope of the Evidence}
The data section states the independent grouping and acquisition scope once, and all results retain that interpretation. Repeated images, seeds, and point coordinates remain correlated. The classification folds cover the available groups but are internal, and the fixed geometry test covers two of those groups. Re-analysis can reveal conditional effects and errors; it cannot create additional subjects or a prospectively untouched test set. No p-value or population confidence interval is therefore computed from the 432 point records as if they were independent tyres.

Mileage categories are proxies, not measured tread depth or roadworthiness. Tread design, background, illumination, or capture conditions may contribute to prediction; the baseline variation motivates that concern but does not identify its cause. The available manual labels lack a blinded repeat-annotation estimate. The point-versus-mask comparison differs in supervision, pretrained representations, and capacity. These limits define the applicable conclusions rather than invalidate the observed paired results.

Video examples expose failure but have no labelled prevalence estimate. Local timing checks do not establish sustained throughput. Target-assisted alignment is software with synthetic tests; a real reference-angle experiment is required for physical accuracy. The article consequently centres on image-space evaluation and does not use alignment screenshots as measurement results.

\subsection{Experiments That Would Add Independent Evidence}
The highest-value next boundary experiment is rotating held-out tyre groups with frozen model settings and fresh training for each split, while retaining a separate development allocation. Existing checkpoints cannot be evaluated on their training tyres and relabelled as cross-validation. A shared-backbone point-versus-mask ablation would separate representation from backbone more directly. Both require additional training; they are not small post-processing analyses and were not executed for this revision.

A blinded second annotation pass on a stratified subset would estimate label repeatability. A tyre-disjoint labelled video set would permit missing-boundary, crossing, point-error, and temporal-consistency rates. Physical claims require depth or wheel-angle references with repeated mounting and runout assessment. These studies address different uncertainties; adding more architectures on the same split would not substitute for them.

\section{Artifact Availability and Reproducibility}
\label{sec:artifacts}
Source code and documentation are available at
\url{https://github.com/Shanmuk4622/tyre-wear-alignment-cv}.
Persistent study artifacts are at
\url{https://huggingface.co/datasets/Shanmuk4622/tyre-wear-study}.
The reporting snapshot is pinned to
\path{22d5a6bc9f953ba3bf2a75919edc7db3193b317b}; the HRNet report to
\path{a92c0f9c5c1b78c6a06e13d51e18722195230658}; and the matched comparison to
\path{bbe586c6f00cf12ae4cac8b2e9cb4f875b272abb}.

The manuscript package includes numerical source hashes and a repeatable revision analysis. The latter checks split membership, reconstructs endpoints from 18 complete epoch histories, repeats explanation-gate decisions, and aggregates paired boundary and fusion records. These are analyses of frozen results, not new training runs. The revision report records every reviewer issue, the implemented response, and remaining evidence needs.

Relevant dense-task versions include Ultralytics 8.4.20, Transformers 4.51.3, segmentation-models-pytorch 0.5.0, and timm 1.0.15. Local integration uses PyTorch 2.5.1 with CUDA 12.1 and OpenCV 4.11.0. Per-phase environments, data lineage, frozen notebooks, and model identities remain necessary for scientific reproduction; a successful paper build alone is not a replication of training.

\section{Conclusion}
The study evaluates three links in a vision-based tyre inspection workflow: checkpoint policy, localisation interventions, and boundary representation. Fold-specific ordinal diagnostics show that pooled selected scores obscure checkpoint and group dependence. Predicted crops and fixed fusion reduce the frozen classifier's mean F1 overall, with the effect concentrated on one fold. Point learning reduces matched horizontal error by 23.75\% relative to mask-derived boundaries and improves 22 of 24 image averages, with complete comparator point coverage. The gain describes the stated pipelines and fixed evaluation cohort. These results support traceable component evaluation and failure visibility; broader transfer and physical measurement remain separate validation tasks.

\section*{Acknowledgment}
No external funding was received for this work. OpenAI Codex~\cite{codex} was used extensively to assist initial drafting and revision of the abstract, Sections I--VIII, appendices, captions, and analysis/build code, including literature organisation, language editing, and checks against stored experimental records. It did not supply experimental observations or independent physical validation. This is an author-review draft; the authors must verify all claims and take responsibility for the final submitted content.

\section*{Declarations for Author Review}
\textit{Proposed contributions (requiring author confirmation):} Sreenivasa Reddy Edara: supervision, methodology review, and writing--review and editing. Shanmukesh Bonala: investigation, software, data curation, formal analysis, visualisation, and writing--original draft. These role assignments are a proposed statement, not a verified record of individual work.

\textit{Competing interests:} author confirmation pending. \textit{Image collection and publication permissions:} author confirmation pending. These fields must be completed before submission.
